\begin{proofbox}
	\textbf{\textcolor{CProof}{证明思路}}

	设 $f$、$g$ 的公共定义域为非空集合 $D$。
	因为两函数各自为奇函数或偶函数，$D$ 关于原点对称。
	对任意 $x\in D$，把 $x$ 换成 $-x$，比较运算结果与原结果。

	\medskip
	\textbf{\textcolor{CProof}{一、加减法}}

	若 $f$、$g$ 都是奇函数，则
	\[
		\begin{aligned}
			(f+g)(-x)
			 & =f(-x)+g(-x) \\
			 & =-f(x)-g(x)  \\
			 & =-(f+g)(x),  \\[2mm]
			(f-g)(-x)
			 & =f(-x)-g(-x) \\
			 & =-f(x)+g(x)  \\
			 & =-(f-g)(x).
		\end{aligned}
	\]
	因此奇函数的和、差仍为奇函数。

	若 $f$、$g$ 都是偶函数，则
	\[
		(f\pm g)(-x)=f(x)\pm g(x)=(f\pm g)(x),
	\]
	因此偶函数的和、差仍为偶函数。

	\medskip
	\textbf{\textcolor{CProof}{二、乘法}}

	分别记 $f(-x)=\varepsilon_f f(x)$、
	$g(-x)=\varepsilon_g g(x)$，其中奇函数对应
	$\varepsilon=-1$，偶函数对应 $\varepsilon=1$。于是
	\[
		\begin{aligned}
			(fg)(-x)
			 & =f(-x)g(-x)                          \\
			 & =\varepsilon_f\varepsilon_g f(x)g(x) \\
			 & =\varepsilon_f\varepsilon_g(fg)(x).
		\end{aligned}
	\]
	所以奇乘奇为偶、奇乘偶为奇、偶乘偶为偶。

	\medskip
	\textbf{\textcolor{CProof}{三、除法}}

	令 $D'=\{x\in D\mid g(x)\ne0\}$，假定 $D'$ 非空。
	因为 $g(-x)=\varepsilon_g g(x)$，$g(x)$ 为零当且仅当
	$g(-x)$ 为零，因此 $D'$ 关于原点对称。对任意 $x\in D'$，
	\[
		\begin{aligned}
			\left(\frac{f}{g}\right)(-x)
			 & =\frac{f(-x)}{g(-x)}        \\
			 & =\frac{\varepsilon_f f(x)}
			{\varepsilon_g g(x)}           \\
			 & =\varepsilon_f\varepsilon_g
			\left(\frac{f}{g}\right)(x).
		\end{aligned}
	\]
	最后一步用到 $\varepsilon_g\in\{-1,1\}$，
	所以 $1/\varepsilon_g=\varepsilon_g$。
	除法由此得到与乘法相同的奇偶规则。
\end{proofbox}